We reproduce the YOLOv1 detector of \citet{redmon2016you} from scratch in
PyTorch, port it from the paper's VOC setting to COCO 2017 (eighty classes),
and train it end-to-end with an ImageNet-pretrained ResNet18 backbone, a
convolutional adapter and the paper's two-layer dense detection head.
The substantive contribution of this work is methodological rather than
numerical. The first four training runs, spanning two backbones and a
deliberately wide sweep of learning rates and gradient-clip thresholds, all
produced exactly $\mathrm{mAP}@0.5 = 0$ on the held-out test split despite
training and validation losses that decreased monotonically and looked, in
isolation, like evidence of a learning model. A code audit identified two
independent implementation bugs latent in the loss and in the coordinate
decoder, each of which is individually sufficient to drive the metric to
zero on COCO while remaining invisible to the pre-existing unit-test suite
(which exercised only the VOC index layout) and to training-curve
inspection. Both bugs were fixed with a handful of lines and verified by
synthetic round-trip and loss-sensitivity tests. The two post-fix runs (V5
and V6) produced the first non-zero detection scores, with the V6 checkpoint
reaching $\mathrm{mAP}@0.5 = 0.0415$ on the test split at a forward-only
latency of $\sim\!5.3$\,ms per $448\!\times\!448$ image on a single
Ada-generation laptop GPU. We do not claim competitiveness with current
detectors; the gap to the paper's $63.4$ on VOC is accounted for by the
harder dataset, a halved training schedule, a smaller backbone and a minimal
augmentation pipeline, each of which we analyse in turn. The takeaway we
hope to leave with the reader is that a converging loss curve is not, on
its own, evidence that the model is learning what the loss is supposed to
measure; cheap quantitative round-trip tests on each coordinate
transformation would have caught both bugs before the first epoch.
